% ---------------------------------------------------------------------
\subsection{Decentralized Uplink}\label{subsec:dbp}
% ---------------------------------------------------------------------

% FONTE: TG1 (sbc-template.tex), Decentralized MIMO Architectures.
The complexity of estimating the vector s grows with the number of antennas and poses a challenge for the MIMO system \cite{jeon2019}. \textcolor{blue}{The Decentralized Baseband Processing (DBP) architecture was introduced in~\cite{decentralized2017, bertilsson2018distributed} to address this issue.} In DBP, the $M$ base-station antennas are partitioned into $C$ clusters with $M_c$ antennas, each cluster with a Distributed Unit (DU) with local radio-frequency circuitry and baseband processing fabric such as custom ASICs, FPGAs, or DSPs \cite{decentralized2017}. 

The uplink system model for a $C$ cluster becomes:

\begin{equation}
    \mathbf{y}_c = \mathbf{H}_c\mathbf{s} + \mathbf{n}_c,
\end{equation}
where $\mathbf{y}_c \in \mathbb{C}^{M_c \times 1}$,
$\mathbf{H}_c \in \mathbb{C}^{M_c \times K}$, and $\mathbf{n}_c$ is the
local noise vector, with ~\cite{jeon2019}:

\begin{equation}
M = \sum_{c=1}^{C} M_c    
\end{equation}

There are two main topologies for the DBP: star, in which all DUs are connected to a central unit (CU), and daisy-chain, in which the DUs are connected in series via unidirectional links and pass intermediate data along  \cite{rodriguezsanchez2020,zhao2025jsac}. This work employs a daisy-chain topology.
% FONTE: zhao2025jsac paper.md:87-89 e :367; rodriguezsanchez2020 L167.
